\documentclass[letterpaper]{article}
\usepackage[utf8]{inputenc}
\usepackage[T1]{fontenc}
\usepackage[activeacute,spanish]{babel}
\usepackage[vmargin=4cm,tmargin=3cm,hmargin=2cm,letterpaper]{geometry}%
\usepackage{helvet}
\usepackage{amsmath,amsfonts,amssymb}
\usepackage{graphicx}
\usepackage{color}
\usepackage{xcolor}
\usepackage{verbatim}
\usepackage{tabls}
\usepackage{lastpage}
\usepackage{fancyhdr}
\usepackage{url}
\usepackage{listings}
\usepackage{tikz}
\usepackage{pgf}
\usepackage{pgffor}
\usepgfmodule{plot}
\usepackage{wrapfig}
\usepackage{ifpdf}
\usepackage{amssymb}
\usepackage{pifont}
\usepackage{epstopdf}
\usepackage{graphicx} % Allows including images
\usepackage{booktabs} % Allows the use of \toprule, \midrule and \bottomrule in tables
\usepackage{tikz}
\usetikzlibrary{arrows,decorations,snakes,backgrounds,fit,calc,through,scopes,positioning,automata,chains,er,fadings,calendar,matrix,mindmap,folding,patterns,petri,plothandlers,plotmarks,shadows,shapes,shapes.arrows,topaths,trees}

\lstset{% general command to set parameter(s)
%   basicstyle=\small,
  % print whole listing small
%   keywordstyle=\color{black}\bfseries\underbar,
  % underlined bold black keywords
%   identifierstyle=,
  % nothing happens
%   commentstyle=\color{white}, % white comments
%   stringstyle=\ttfamily,
  % typewriter type for strings
  showstringspaces=false}
  % no special string spaces

\pagestyle{fancy}
\color{black}
\fancyhead{}
\renewcommand{\headrule}{\hrule\vspace*{0.5mm}\rule{\linewidth}{0.8mm}}
\renewcommand{\familydefault}{\sfdefault}

\graphicspath{{./images/}}
\lhead{\includegraphics[width=2cm]{logoucr.png}}
\rhead{\includegraphics[width=3cm]{eie-text-gray-6x3cm.png}}
\chead{UNIVERSIDAD DE COSTA RICA\\FACULTAD DE INGENIERÍA\\ESCUELA DE INGENIERÍA ELÉCTRICA\\\textbf{ESTRUCTURAS ABSTRACTAS DE DATOS Y\\ ALGORITMOS PARA INGENIERÍA}\\IE-0217\\I CICLO 2014\\PROPUESTA DE I PROYECTO LIBRERÍA C++}

\lfoot{}%
\cfoot{}%
%\cfoot{\thepage\ de \pageref{LastPage}}%
\rfoot{}%

%%%%%%%%%%%%%%%%%%%%%%%%%%%%%%%%%%%%%%%%%%%%%%%%%%%%%%%%%%%%%%%%%%%%%%%%%%%%%%%%%%%%%%%%%%%%%%%%%%%%%%%%%%%%%%%
\newcommand{\uic}{blue} %user-input color
%%%%%%%%%%%%%%%%%%%%%%%%%%%%%%%%%%%%%%%%%%%%%%%%%%%%%%%%%%%%%%%%%%%%%%%%%%%%%%%%%%%%%%%%%%%%%%%%%%%%%%%%%%%%%%%%%%
\newcommand{\uim}{\_\_} %user-input marker
%%%%%%%%%%%%%%%%%%%%%%%%%%%%%%%%%%%%%%%%%%%%%%%%%%%%%%%%%%%%%%%%%%%%%%%%%%%%%%%%%%%%%%%%%%%%%%%%%%%%%%%%%%%%%%%%%%
\newcommand{\userinput}[1]{\textcolor{\uic}{\uim#1\uim}}


%%%%%%%%%%%%%%%%%%%%%%%%%%%%%%%%%%%%%%%%%%%%%%%%%%%%%%%%%%%%%%%%%%%%%%%%%%%%%%%%%%%%%%%%%%%%%%%%%%%%%%%%%%%%%%%%%%
\begin{document}\vspace*{2cm}
%%%%%%%%%%%%%%%%%%%%%%%%%%%%%%%%%%%%%%%%%%%%%%%%%%%%%%%%%%%%%%%%%%%%%%%%%%%%%%%%%%%%%%%%%%%%%%%%%%%%%%%%%%%%%%%%%%

%%%%%%%%%%%%%%%%%%%%%%%%%%%%%%%%%%%%%%%%%%%%%%%%%%%%%%%%%%%%%%%%%%%%%%%%%%%%%%%%%%%%%%%%%%%%%%%%%%%%%%%%%%%%%%%%%%
\begin{center}
\Huge
\textbf{Kinect Data Library}
\vspace*{1cm}
\end{center}

\noindent
\small\baselineskip=14pt
\textbf{Estudiantes:} \\
\text{David Pérez Bolaños - B04769}\\
\text{Andrey Pérez - BXXXXX}\\
\text{Andrés Sánchez López - B26214}\\


%%%%%%%%%%%%%%%%%%%%%%%%%%%%%%%%%%%%%%%%%%%%%%%%%%%%%%%%%%%%%%%%%%%%%%%%%%%%%%%%%%%%%%%%%%%%%%%%%%%%%%%%%%%%%%%%%%
\section{Introducción}
\hyphenation{fun-cio-na-li-dad}
La creación de librerías en lenguajes de programación ayuda a generar una interfaz bien definida para una cierta funcionalidad en específico, estas sirven para separar por módulos un programa, y así generar un código más claro y ordenado.\\

Una librería es un conjunto de funciones para desarrollar software, por lo general no son programas, pero si son utilizadas por los programas para poder funcionar de forma correcta; el desarrollo de librerías sirve como apoyo para los programadores a tener más facilidades de implementación 
en sus programas y a contar con más recursos para realizar sus proyectos.\\

En este caso, queremos implementar una librería sobre manipulación y análisis de objetos tridimensionales
en el lenguaje C++.\\

La idea  de esta librería es crear métodos o funciones para la manipulación y representación 
de cualquier objeto real, de manera virtual; esta técnica de pasar de un objeto real a representarlo de manera virtual, se utiliza mucho en videojuegos, efectos especiales, medicina, simuladores, entre otros.\\

 Para nuestro caso utilizaremos para representar los objetos de manera virtual, la técnica de malla de triángulos en 3D, en dónde primero debemos representar por medio de una nube de puntos el objeto y luego formar los triangulos con cada uno de esos puntos; para ello utilizaremos el kinect, obteniendo así los datos y luego realizar los algoritmos necesarios para la ejecución de las funciones de comparación, simetría, impresión de vistas y cortes sobre los datos mismos, generando así la librería. 

%%%%%%%%%%%%%%%%%%%%%%%%%%%%%%%%%%%%%%%%%%%%%%%%%%%%%%%%%%%%%%%%%%%%%%%%%%%%%%%%%%%%%%%%%%%%%%%%%%%%%%%%%%%%%%%%%%
\section{Objetivos}

\subsection{Objetivo General}


El objetivo general consiste en la creación de una librería en C++ para la manipulación de objetos tridimensionales.\\


\subsection{Objetivos Específicos}

Los objetivos específicos son:\\

\begin{enumerate}
\item Ampliar el conocimiento del lenguaje C++ a través de la creación de una librería en este lenguaje. 
\item Desarrollar una función de comparación entre las posiciones y las trayectorias de figuras tridimensionales en la librería.
\item Ajustar una función de impresión de las coordenadas de posición de una figura tridimensional en código C++ en la librería. 
\item Realizar una función de cortes en las coordenadas de posición de una figura tridimensional en la librería.
\item Implementar una función que consiga un conjunto de coordenadas de posición que sean simétricas a otro conjunto de coordenadas dado.
\item Obtener los datos para la revisión de las funciones de la librería utilizando un kinect.
\end{enumerate}

%%%%%%%%%%%%%%%%%%%%%%%%%%%%%%%%%%%%%%%%%%%%%%%%%%%%%%%%%%%%%%%%%%%%%%%%%%%%%%%%%%%%%%%%%%%%%%%%%%%%%%%%%%%%%%%%%%
\section{Metodología}


Se realizará la investigación sobre librerías que se necesiten para la manipulación de objetos tridimensionales. Previamente, se revisarán las 
librerías encontradas, para determinar la necesidad de nuestro proyecto. 
Una vez hecho esto, se buscará implementar una librería en el lenguaje C++.\\

Lo que se realizará primero, es una investigación sobre que funciones presentan las librerías encontradas en lenguaje C++ y estudiar el código de estas,
esto con el fin de entender la utilización de los datos en estos códigos, todo esto, en caso de haber encontrado una librería en lenguaje C++ con relación a este tema.
Para lograr esto, se leerá toda la bibliografía necesaria para poder entender de la mejor manera que es lo que se necesita exactamente. Una vez
entendido algunas de las diferentes implementaciones ya existentes, se pretende crear una librería en lenguaje C++ con diferentes funciones, entre ellas,
comparación e impresión de figuras tridimensionales, así como cortes y simetrías de la misma. \\

Será necesario también, la utilización de un
kinect para obtener datos de diferentes figuras tridimensionales de tal modo que se pueda dar una revisión de la librería. Finalmente, se buscará ajustar tal librería,
de modo que sea fácil para el usuario la utilización de la misma, para que sea utilizada como una solución de las necesidades
existentes sobre objetos tridimensionales.\\

%%%%%%%%%%%%%%%%%%%%%%%%%%%%%%%%%%%%%%%%%%%%%%%%%%%%%%%%%%%%%%%%%%%%%%%%%%%%%%%%%%%%%%%%%%%%%%%%%%%%%%%%%%%%%%%%%%

\section{Cronograma}
\begin{center}
\begin{tabular}{l l   @{\hspace{1cm}}p{10cm}}
\cline{3-3}

\toprule
\textbf{Semana} & \textbf{Fechas} & \textbf{Actividad} \\
\midrule
1 & 13 de abril a 19 de abril & Entrega de la propuesta - Instalación del
software para uso del kinect - Estudio del software para uso del kinect -
Estudio del lenguaje C++ \\
2 & 20 de abril a 26 de abril & Desarrollo, en lenguaje C++, de la función
de comparación e impresión. \\
3 & 27 de abril a 3 de mayo & Desarrollo, en lenguaje C++, de la función
de cortes y simetrías. \\
4 & 4 de mayo a 10 de mayo & Pruebas de las 4 funciones utilizando datos
provenientes del kinect. \\
5 & 11 de mayo a 15 de mayo & Realización de la presentación e informe
escrito del proyecto. \\
\bottomrule
\end{tabular}
\end{center}



\section{Referencias}

\begin{enumerate}

\item Richard, J. Computer Science Division. University of California at Berkeley. Triangle. A Two-Dimensional Quality Mesh Generator and 
Delaunay Triangulator. Encontrado el 13 de abril del 2014 en: http://www.cs.cmu.edu/~quake/triangle.html
\item Escenografía Intermedial. Nuevos medios y tecnologías afines a la escena. (15 de mayo del 2012).
Nube de puntos (Point Cloud) con Kinect. Encontrado el 13 de abril del 2014 en: http://escenografiaaumentada.wordpress.com/2012/05/15/148/
\item OPENKINECT. Encontrado el 13 de abril del 2014 en: http://openkinect.org/wiki/Main\_Page

\end{enumerate}
	
\end{document}

